\section{Conclusion}
\label{sec:conclusion}

ACDB is a measurement instrument first and an empirical study second. Its design choices — two retained truth objects driving a layered scoring contract, a random-floor-calibrated ladder, paired accepted-seed manifests across methods, and hybrid ablations that cut the pipeline at the joints SCM theory already identifies — are what make differences between methods attributable rather than merely measured. The §\ref{sec:results} patterns illustrate the point: that statistical-tool access does not reliably improve directed recovery, that supplying the PC-derived CPDAG closes most of the gap for weaker models but barely moves \texttt{GPT-5.5}, and that end-to-end LLM policies degrade faster than \texttt{pc\_greedy} with graph scale, are interpretable as separable claims about observational recovery, interventional orientation, and end-to-end integration only because the instrument separates those stages. Without that separation, the same evidence would collapse into a single ranking and the same models would be called good or bad at "causal reasoning" with no further structure.

The fully observed linear--Gaussian regime with hard single-node interventions is the precondition for that separation. It is the setting in which every score layer has an exact theoretical referent: CPDAG identifiability is unambiguous, the resolving-intervention set is well-defined, and the structure-blind random floor admits a closed form. The substantive limitations this entails for non-linear, partially observed, or semantic substrates are stated explicitly in §\ref{sec:limitations}; widening the instrument to those regimes is scientifically valuable but reintroduces exactly the ambiguity the present version was built to remove. ACDB is therefore a base measurement layer, designed so that future extensions inherit its scoring discipline rather than replace it.

Treated as such, the instrument enables questions this paper does not yet answer. Whether a new model's gains on causal reasoning come from observational pattern matching, from intervention choice, or from sequential evidence integration becomes a question the layered scores and hybrid ablations can address directly. Whether tool augmentation or fine-tuning moves the residual orientation gap, or only the observational layer, becomes a controlled comparison rather than a debate. Whether the LLM-vs-classical degradation gap closes at larger graph scales for any model family becomes a measurable curve. ACDB is built so that the conclusions it supports remain legible as both the scientific environments and the agents become more capable.
